\subsection{\texorpdfstring{第~\ref{sec:hawking-singularity-theorem}~节}{Hawking 奇性定理}}\label{subsec:solutions-hawking-singularity-theorem}

\begin{description}[font=\normalfont,leftmargin=2em,labelwidth=1.5em]
\RGsolution{ex:clifton-pohl-torus-incompleteness}{clifton-pohl-torus-incompleteness}
$(\overline M,\overline g)$ 上仅有的非零 Christoffel 符号为
\[
  \Gamma_{vv}^v=-\frac{2v}{v^2+\mathit{v}^2}
  \qquad\text{且}\qquad
  \Gamma_{\mathit{v}\mathit{v}}^{\mathit{v}}
    =-\frac{2\mathit{v}}{v^2+\mathit{v}^2}.
\]
因此测地线方程为
\[
  \ddot v-\frac{2v}{v^2+\mathit{v}^2}\dot v^2=0
\]
和
\[
  \ddot{\mathit{v}}-\frac{2\mathit{v}}{v^2+\mathit{v}^2}\dot{\mathit{v}}^2=0.
\]
例如, 取 $\mathit{v}\equiv0$, 得
\[
  \frac{\ddot v}{\dot v}=2\frac{\dot v}{v},
\]
于是
\[
  \dot v=-av^2,
\]
从而
\[
  v(t)=\frac{1}{at+b}
\]
对某些常数 $a,b\in\mathbb{R}$ 成立. 因此, 曲线
\[
  c(t):=\left(\frac{1}{at+b},0\right)
\]
是经过点 $c(0)=(\frac{1}{b},0)$ 的零模测地线. 例如, 若 $b>0$ 且 $a<0$, 则这条测地线在 $t<-\frac{b}{a}$ 上有定义. 这条测地线的像为
\[
  \{(v,0)\mid v>0\},
\]
所以这条曲线在 $\overline M$ 中不可延拓. 因此 $\overline M$ (从而 $M$) 不是测地完备的.

\RGsolution{ex:locally-finite-refinements-by-simple-neighborhoods}{locally-finite-refinements-by-simple-neighborhoods}
显然, 简单邻域构成 $M$ 的拓扑的一组基, 因而对于每个开覆盖 $\{V_\alpha\}_{\alpha\in A}$, 都存在由简单邻域构成的加细 $\{U_\beta\}_{\beta\in B}$, 即
\[
  \bigcup_{\alpha\in A}V_\alpha=\bigcup_{\beta\in B}U_\beta
\]
且对每个 $\beta\in B$ 都存在 $\alpha\in A$ 使 $U_\beta\subset V_\alpha$.

先假定 $V:=\bigcup_{\alpha\in A}V_\alpha$ 紧. 则存在有限子覆盖 $\{U_{\beta_i}\}_{i=1}^k$, 使
\[
  V=\bigcup_{i=1}^k U_{\beta_i}.
\]
显然, $\{U_{\beta_i}\}_{i=1}^k$ 是开覆盖 $\{V_\alpha\}_{\alpha\in A}$ 的一个由简单邻域构成的可数局部有限加细.

若 $V$ 不紧, 可利用紧穷竭, 即 $V$ 的一列紧子集 $\{K_i\}_{i\in\mathbb{N}}$, 满足 $K_i\subset K_{i+1}$ 且 $M=\bigcup_{i=1}^\infty K_i$ (参见第~\ref{chap:differential-forms}~章的 Remark~\ref{rem:compact-exhaustions-of-manifolds}). 族 $\{U_\beta\}_{\beta\in B}$ 是 $K_1$ 的一个覆盖, 故可取 $K_1$ 的有限子覆盖
\[
  \{U_{\beta_1},\ldots,U_{\beta_{k_1}}\}.
\]
通过归纳, 得到有限个邻域
\[
  \{U_{\beta_1^i},\ldots,U_{\beta_{k_i}^i}\}
\]
覆盖 $K_i\setminus\operatorname{int}K_{i-1}$ (一个紧集). 注意, 必要时可改取较小的简单邻域, 使
\[
  \bigcup_{j=1}^{k_i}U_{\beta_j^i}\subset\operatorname{int}K_{i+1}\setminus K_{i-2},
\]
从而可数覆盖 $\{U_{\beta_j^i}\}_{i\in\mathbb{N},\,1\leq j\leq k_i}$ 局部有限. 结论随之得出.

\RGsolution{ex:prove-closed-causal-futures-and-compact-diamonds}{prove-closed-causal-futures-and-compact-diamonds}
先注意, 用与 Proposition~\ref{prop:compact-past-intersect-future-domain-of-dependence} 类似的论证可证 $D^-(S)\cap J^+(p)$ 紧.
\begin{enumerate}[label=(\roman*),ref=\roman*,leftmargin=*]
\item 考虑 $J^+(p)$ 中收敛于 $q\in M$ 的点列 $q_n$, 并取满足 $t(S)>t(q)$ 的 Cauchy 曲面 $S$ (其中 $t:M\to\mathbb{R}$ 是 $M$ 的全局时间函数), 使得 $q\in D^-(S)$. 则对于充分大的 $n\in\mathbb{N}$, 有 $q_n\in A:=D^-(S)\cap J^+(p)$. 由于 $A$ 紧, 可知 $q\in A\subset J^+(p)$.
\item 首先注意, (i) 对 $J^-(q)$ 同样成立, 即对任何 $q\in M$, $J^-(q)$ 都是闭集. 因而 $B:=J^+(p)\cap J^-(q)$ 为闭集. 取满足 $t(S)=t(q)$ 的 Cauchy 曲面 $S$, 则 $B$ 是紧集 $D^-(S)\cap J^+(p)$ 的闭子集, 因此其本身也是紧集.
\end{enumerate}

\RGsolution{ex:hausdorff-closure-of-timelike-curves}{hausdorff-closure-of-timelike-curves}
设 $\gamma\in C(S,p)$. 则存在序列 $\gamma_n\in T(S,p)$, 使 $\gamma_n\to\gamma$.

先证明当 $0\leq a\leq t(p)$ 时, $\gamma$ 与每个水平集 $S_a:=t^{-1}(a)$ 相交. 事实上, 若 $\gamma\cap S_a=\varnothing$, 那么由于 $\gamma$ 紧而 $S_a$ 闭, 两者之间的距离大于某个 $\varepsilon>0$. 但每条 $\gamma_n$ 都与 $S_a$ 相交, 因而有 $d_H(\gamma,\gamma_n)>\varepsilon$, 矛盾. 类似地可证, 当 $a<0$ 或 $a>t(p)$ 时, $\gamma$ 不可能与 $S_a$ 相交.

不难验证, 由 $\pi_a(c)=c\cap S_a$ 给出的映射 $\pi_a:C(S,p)\to C(S_a)$ 连续 (这里 $C(S_a)$ 是 $S_a$ 的所有紧子集构成的带 Hausdorff 度量的空间). 因此 $\pi_a(\gamma_n)\to\pi_a(\gamma)$. 由于每个 $\pi_a(\gamma_n)$ 都是一个点, $\pi_a(\gamma)$ 也是一个点, 从而可将 $\gamma$ 视为映射 $\gamma:[0,t(p)]\to M$.

为证明 $\gamma$ 是因果曲线, 注意当 $0\leq a<b\leq t(p)$ 时, 有 $\gamma_n(a)\to\gamma(a)$ 且 $\gamma_n(b)\to\gamma(b)$, 其中 $\gamma_n(b)\in I^+(\gamma_n(a))$. 若 $q\in I^-(\gamma(a))$, 则当 $n$ 充分大时有 $\gamma_n(a)\in I^+(q)$, 进而 $\gamma_n(b)\in I^+(q)$. 于是 $\gamma(b)\in\overline{I^+(q)}=J^+(q)$ (回忆 $J^+(q)\subset\overline{I^+(q)}$ 是闭集), 从而 $q\in J^-(\gamma(b))$. 取点列 $q_n\in I^-(\gamma(a))$ 使 $q_n\to\gamma(a)$, 可得 $\gamma(a)\in J^-(\gamma(b))$.

最后, $\gamma$ 必连续: 若 $a_n\in[0,t(p)]$ 为极限是 $a\in[0,t(p)]$ 的递增序列, 则 $\gamma(a_n)\to\gamma(a)$. 否则存在子序列 $b_n$, 使 $\gamma(b_n)\to q\in S_a$, 并且当 $n$ 充分大时, 有 $\gamma(b_n)\notin J^-(\gamma(a))$.

\RGsolution{ex:myers-compactness-under-positive-ricci-curvature}{myers-compactness-under-positive-ricci-curvature}
取 $p\in M$ 及法球面 $S=S_\delta(p)$. 由于 $(M,g)$ 完备, 可利用指向外的单位法向定义映射 $\exp:\mathbb{R}\times S\to M$, 且 $M=B_\delta(p)\cup\exp(\mathbb{R}\times S)$. 因为 $S$ 紧, 在 $S$ 上有 $\theta\leq\theta_0$, 其中 $\theta_0>0$. 与 Proposition~\ref{prop:timelike-focusing-under-strong-energy-condition} 的证明完全相同的计算表明
\[
  \frac{\partial\theta}{\partial t}+\frac{1}{n-1}\theta^2\leq-\varepsilon,
\]
其中 $n=\dim M$. 特别地,
\[
  \frac{\partial\theta}{\partial t}\leq-\varepsilon,
\]
所以沿任意测地线, 存在某个 $t\leq\frac{1+\theta_0}{\varepsilon}$ 使 $\theta\leq-1$ (注意直线 $\theta=-\varepsilon t+\theta_0$ 与 $\theta=-1$ 在 $t=\frac{1+\theta_0}{\varepsilon}$ 处相交). 自该点起有
\[
  \frac{\partial\theta}{\partial t}+\frac{1}{n-1}\theta^2\leq0
  \quad\Longrightarrow\quad
  \frac{1}{\theta}\geq\frac{t}{n-1}-1.
\]
因此每条与 $S$ 正交且指向外的测地线在不超过
\[
  \frac{1+\theta_0}{\varepsilon}+n-1
\]
的弧长处到达一个共轭点, 在那里不再是最短的. 因而 $M$ 中所有点到 $p$ 的距离都小于
\[
  \delta+\frac{1+\theta_0}{\varepsilon}+n-1,
\]
这意味着 $M$ 有界, 从而紧.

注意, 在 Riemann 情形下并无矛盾: 一条测地线上越过共轭点的各点, 都能从 $S$ 沿一条更短的 (而非 Lorentz 情形下的更长的) 测地线到达, 因而那条测地线本身不一定已经到达一个共轭点. 所以, 奇性定理的证明不适用于 Riemann 几何.

\RGsolution{ex:limits-of-hawking-singularity-theorem}{limits-of-hawking-singularity-theorem}
\begin{enumerate}[label=(\alph*),ref=\alph*,leftmargin=*]
\item Minkowski 时空不包含膨胀量满足 $\theta\leq\theta_0<0$ 的 Cauchy 超曲面.
\item Einstein 宇宙不包含膨胀量满足 $\theta\leq\theta_0<0$ 的 Cauchy 超曲面.
\item de Sitter 宇宙不满足强能量条件 [参见 Exercise~\ref{ex:hawking-singularity-theorem}(\ref{ex:strong-energy-condition-with-cosmological-constant})].
\item 二维反 de Sitter 时空不是整体双曲的; 但应注意, 该时空的任何整体双曲开子集都不是测地完备的.
\end{enumerate}
\end{description}
